% GENERATED FILE. Edit canonical lessons, then run npm run generate:books.
% canonical-source-hash: fnv1a64:17b0209c070c8de8
% canonical-lessons: FR-C431-reclamation, FR-C431-reclamer, FR-C431-rouspeter, FR-C431-deranger, FR-C431-mecontent, FR-R431-revision-424-431

\chapter{On n'est pas content}
\label{ch:fr-on-n-est-pas-content}
\hlchaptermodality{431}

\begin{chapteropening}
\textbf{By the end of this chapter:} \emph{I can complain: the complaint, demanding, grumbling, disturbing someone and being unhappy.}

Complete the last lesson of chapter 431: I can complain: the complaint, demanding, grumbling, disturbing someone and being unhappy.
\end{chapteropening}

\section[la réclamation]{la réclamation — the complaint, the claim}

\label{lesson:FR-C431-reclamation}



\begin{quote}
Before the new one: say the French for the shutter, then the French for the blind.
\end{quote}

\subsection*{You'll want to know: la réclamation}
\textbf{la réclamation} — \textquotedblleft{}the complaint, the claim\textquotedblright{}.

\begin{quote}
Nous avons fait une réclamation à la direction de l'hôtel.
\end{quote}

\begin{quote}
\emph{We made a complaint to the hotel management.}
\end{quote}

Formal; \emph{le service des réclamations} — the complaints department.

\begin{cousinweb}
\emph{réclamation} is a learned borrowing from Latin \emph{reclāmātiō} 'a crying out against' (EtymDB), from \emph{clāmāre} 'to shout'.
\end{cousinweb}

\begin{grammarlens}[title={What you've built}]
One more word for telling, explaining and describing at B1.
\end{grammarlens}

\subsection*{Your turn}
\begin{itemize}
\raggedright
  \item \emph{Say it:} \emph{la réclamation}
  \item \emph{Say it:} \emph{la réclamation}, once more
  \item \emph{Say it:} a sentence of your own with \emph{la réclamation}
  \item \emph{Recall:} say the French for the shutter, then the French for the blind, then say \emph{la réclamation} again
\end{itemize}

\begin{tcolorbox}[breakable,colback=teal!4,colframe=teal!35!black,title={Before you move on}]
What does la réclamation mean? (\textquotedblleft{}The complaint, the claim\textquotedblright{}.) Say it once more.
\end{tcolorbox}
\section[réclamer]{réclamer — to ask for, to demand; to claim}

\label{lesson:FR-C431-reclamer}



\begin{quote}
Before the new one: say the French for the blind, then the French for the complaint.
\end{quote}

\subsection*{You'll want to know: réclamer}
\textbf{réclamer} — \textquotedblleft{}to ask for, to demand; to claim\textquotedblright{}.

\begin{quote}
Il a réclamé une autre chambre, plus calme.
\end{quote}

\begin{quote}
\emph{He asked for another, quieter room.}
\end{quote}

\emph{Réclamer son argent} — to demand one's money back.

\begin{cousinweb}
\emph{réclamer} comes from Old French \emph{reclamer}, from Latin \emph{reclāmāre} 'to cry out against', built on \emph{clāmāre} 'to shout' (EtymDB).
\end{cousinweb}

\begin{grammarlens}[title={What you've built}]
One more word for telling, explaining and describing at B1.
\end{grammarlens}

\subsection*{Your turn}
\begin{itemize}
\raggedright
  \item \emph{Say it:} \emph{réclamer}
  \item \emph{Say it:} \emph{réclamer}, once more
  \item \emph{Say it:} a sentence of your own with \emph{réclamer}
  \item \emph{Recall:} say the French for the blind, then the French for the complaint, then say \emph{réclamer} again
\end{itemize}

\begin{tcolorbox}[breakable,colback=teal!4,colframe=teal!35!black,title={Before you move on}]
What does réclamer mean? (\textquotedblleft{}To ask for, to demand; to claim\textquotedblright{}.) Say it once more.
\end{tcolorbox}
\section[rouspéter]{rouspéter — to grumble, to moan}

\label{lesson:FR-C431-rouspeter}



\begin{quote}
Before the new one: say the French for the complaint, then the French for to ask for.
\end{quote}

\subsection*{You'll want to know: rouspéter}
\textbf{rouspéter} — \textquotedblleft{}to grumble, to moan\textquotedblright{}.

\begin{quote}
Mon oncle rouspète toujours contre le service.
\end{quote}

\begin{quote}
\emph{My uncle is forever grumbling about the service.}
\end{quote}

Informal. \emph{Il rouspète} — he moans.

\begin{cousinweb}
The origin of \emph{rouspéter} is not settled; it is an old slang word. EtymDB lists it without a source.
\end{cousinweb}

\begin{grammarlens}[title={What you've built}]
One more word for telling, explaining and describing at B1.
\end{grammarlens}

\subsection*{Your turn}
\begin{itemize}
\raggedright
  \item \emph{Say it:} \emph{rouspéter}
  \item \emph{Say it:} \emph{rouspéter}, once more
  \item \emph{Say it:} a sentence of your own with \emph{rouspéter}
  \item \emph{Recall:} say the French for the complaint, then the French for to ask for, then say \emph{rouspéter} again
\end{itemize}

\begin{tcolorbox}[breakable,colback=teal!4,colframe=teal!35!black,title={Before you move on}]
What does rouspéter mean? (\textquotedblleft{}To grumble, to moan\textquotedblright{}.) Say it once more.
\end{tcolorbox}
\section[déranger]{déranger — to disturb, to bother}

\label{lesson:FR-C431-deranger}



\begin{quote}
Before the new one: say the French for to ask for, then the French for to grumble.
\end{quote}

\subsection*{You'll want to know: déranger}
\textbf{déranger} — \textquotedblleft{}to disturb, to bother\textquotedblright{}.

\begin{quote}
La musique des voisins nous a dérangés toute la nuit.
\end{quote}

\begin{quote}
\emph{The neighbours' music disturbed us all night.}
\end{quote}

\emph{Ne pas déranger} — do not disturb, on a hotel door.

\begin{cousinweb}
\emph{déranger} comes from Old French \emph{desrengier} 'to throw out of line', from \emph{des-} and \emph{rengier} 'to put in a row', built on \emph{reng} 'row' (EtymDB).
\end{cousinweb}

\begin{grammarlens}[title={What you've built}]
One more word for telling, explaining and describing at B1.
\end{grammarlens}

\subsection*{Your turn}
\begin{itemize}
\raggedright
  \item \emph{Say it:} \emph{déranger}
  \item \emph{Say it:} \emph{déranger}, once more
  \item \emph{Say it:} a sentence of your own with \emph{déranger}
  \item \emph{Recall:} say the French for to ask for, then the French for to grumble, then say \emph{déranger} again
\end{itemize}

\begin{tcolorbox}[breakable,colback=teal!4,colframe=teal!35!black,title={Before you move on}]
What does déranger mean? (\textquotedblleft{}To disturb, to bother\textquotedblright{}.) Say it once more.
\end{tcolorbox}
\section[mécontent]{mécontent — unhappy, dissatisfied}

\label{lesson:FR-C431-mecontent}



\begin{quote}
Before the new one: say the French for to grumble, then the French for to disturb.
\end{quote}

\subsection*{You'll want to know: mécontent}
\textbf{mécontent} — \textquotedblleft{}unhappy, dissatisfied\textquotedblright{}.

\begin{quote}
Les clients sont mécontents du petit-déjeuner.
\end{quote}

\begin{quote}
\emph{The guests are unhappy with the breakfast.}
\end{quote}

\emph{Mécontent de} — unhappy with; \emph{le mécontentement} — discontent.

\begin{cousinweb}
\emph{mécontent} is built from \emph{mé-} 'badly, not' and \emph{content} 'glad', from Latin \emph{contentus} 'satisfied', from \emph{continēre}, built on \emph{tenēre} 'to hold' (Lewis \& Short). EtymDB does not list it.
\end{cousinweb}

\begin{grammarlens}[title={What you've built}]
One more word for telling, explaining and describing at B1.
\end{grammarlens}

\subsection*{Your turn}
\begin{itemize}
\raggedright
  \item \emph{Say it:} \emph{mécontent}
  \item \emph{Say it:} \emph{mécontent}, once more
  \item \emph{Say it:} a sentence of your own with \emph{mécontent}
  \item \emph{Recall:} say the French for to grumble, then the French for to disturb, then say \emph{mécontent} again
\end{itemize}

\begin{tcolorbox}[breakable,colback=teal!4,colframe=teal!35!black,title={Before you move on}]
What does mécontent mean? (\textquotedblleft{}Unhappy, dissatisfied\textquotedblright{}.) Say it once more.
\end{tcolorbox}
\section[(dialogue)]{Retour sur les chapitres 424–431}

\label{lesson:FR-R431-revision-424-431}



\begin{quote}
Say the French for the inn and the holiday cottage.
\end{quote}

\subsection*{Your turn}
Say these aloud:

\begin{itemize}
\raggedright
  \item the inn, the holiday cottage, the overnight stay, to lodge, to stay for a while
  \item the campsite, the tent, the sleeping bag, to camp, to bivouac
  \item the reception desk, the receptionist, the storey, the stay, to unpack
  \item the air conditioning, the heating, the radiator, the pillowcase, the bathrobe
  \item the leak, the tap, the washbasin, to drip, to unblock
  \item the light bulb, the socket, the minibar, the fitted carpet, the stain
  \item the bedbug, the cockroach, the mosquito net, the shutter, the blind
  \item the complaint, to ask for, to grumble, to disturb, unhappy
\end{itemize}

\begin{tcolorbox}[breakable,colback=teal!4,colframe=teal!35!black,title={Before you move on}]
The inn? (\textbf{l'auberge}.) The holiday cottage? (\textbf{le gîte}.) The overnight stay? (\textbf{la nuitée}.) To lodge? (\textbf{loger}.) To stay for a while? (\textbf{séjourner}.) The campsite? (\textbf{le camping}.) The tent? (\textbf{la tente}.) The sleeping bag? (\textbf{le sac de couchage}.) To camp? (\textbf{camper}.) To bivouac? (\textbf{bivouaquer}.) The reception desk? (\textbf{la réception}.) The receptionist? (\textbf{le réceptionniste}.) The storey? (\textbf{l'étage}.) The stay? (\textbf{le séjour}.) To unpack? (\textbf{déballer}.) The air conditioning? (\textbf{la climatisation}.) The heating? (\textbf{le chauffage}.) The radiator? (\textbf{le radiateur}.) The pillowcase? (\textbf{la taie}.) The bathrobe? (\textbf{le peignoir}.) The leak? (\textbf{la fuite}.) The tap? (\textbf{le robinet}.) The washbasin? (\textbf{le lavabo}.) To drip? (\textbf{goutter}.) To unblock? (\textbf{déboucher}.) The light bulb? (\textbf{l'ampoule}.) The socket? (\textbf{la prise}.) The minibar? (\textbf{le minibar}.) The fitted carpet? (\textbf{la moquette}.) The stain? (\textbf{la tache}.) The bedbug? (\textbf{la punaise}.) The cockroach? (\textbf{le cafard}.) The mosquito net? (\textbf{la moustiquaire}.) The shutter? (\textbf{le volet}.) The blind? (\textbf{le store}.) The complaint? (\textbf{la réclamation}.) To ask for? (\textbf{réclamer}.) To grumble? (\textbf{rouspéter}.) To disturb? (\textbf{déranger}.) Unhappy? (\textbf{mécontent}.)
\end{tcolorbox}
